\section{总结与延伸}

\subsection{核心结论}

\begin{importantbox}{八条可迁移结论}
\begin{enumerate}
  \item AI coding 是 context、model、apply、validation 的闭环，各层需要不同 SLO。
  \item Blast radius 由 runtime/state boundary 决定，不由 service 数量决定。
  \item Merkle-style sync、content cache 和 retrieval evaluation 把全量工作变为可测 delta work。
  \item Metadata、jobs、source objects 与 derived index 需要不同 source of truth 和 rebuild path。
  \item Retry、cron、MVCC bloat 和 rebuild 会形成事故正反馈，恢复时先停止 producer。
  \item Migration 的隐藏成本是 re-embedding、dual validation、cache warming 和临时 capacity。
  \item Object storage 解耦耐久与 compute；global inference 仍需 cache、region、provider eval 与 fallback。
  \item Security、abuse、tool approval 与 incident response 都是产品持续性的基础设施。
\end{enumerate}
\end{importantbox}

\subsection{实践作业：设计一个 AI coding backend}

为一个支持 autocomplete、repo chat 与 multi-file agent 的系统完成以下设计：

\begin{enumerate}
  \item 定义 indexing、inference、apply 和 data plane 的 SLO 与 blast radius；
  \item 设计 Merkle-style sync、chunk/version key、embedding cache 和 deletion semantics；
  \item 列出 metadata、job、source object、vector segment 的 source of truth、rebuild path 与 migration validation；
  \item 为 retry storm、database disk pressure、object-storage cache miss 和 provider rate limit 写 degraded mode；
  \item 建立 self-hosted/frontier routing、offline/online eval、privacy、quota、tool approval 与 incident command。
\end{enumerate}

\begin{knowledgebox}{验收标准}
高质量方案不会只列“vector DB + LLM + queue”。它应说明可重建状态、retry 放大、dual validation、model version telemetry，以及代码与 terminal 权限怎样被理解和撤销。
\end{knowledgebox}

\subsection{拓展阅读}

建议依次阅读 Cursor 的 \href{https://cursor.com/blog/secure-codebase-indexing}{secure indexing}、\href{https://cursor.com/blog/instant-apply}{Fast Apply}、\href{https://cursor.com/blog/cursorbench}{CursorBench}、\href{https://cursor.com/docs/enterprise/privacy-and-data-governance}{privacy} 与 \href{https://cursor.com/docs/agent/security}{agent security}；再读 turbopuffer 的 \href{https://turbopuffer.com/blog/turbopuffer}{object-storage search} 和 \href{https://turbopuffer.com/blog/continuous-recall}{continuous recall}；最后对照 PostgreSQL \href{https://www.postgresql.org/docs/current/routine-vacuuming.html}{vacuum} 与 Amazon S3 \href{https://docs.aws.amazon.com/AmazonS3/latest/userguide/Welcome.html}{object/consistency model}。

\end{document}
